\section{Actual realization and complete stopped paths}
\label{sec:growth-realization}
\label{sec:physical-paths}

The comparison in Section~\ref{sec:graph} controls actual touchdown
states and the first physical section. We now recover the whole
motion between those events. On fixed positive input bands we pair
all contact and flight prefixes. A continuous scan of one unchanged
contact index realizes its stopping phase; matching finitely many
last returns and controlling the entire earlier tail then gives the
complete path.

\subsection{Ordered comparison on every positive-scale compact set}

The realization argument scans $[s/2,2s]$, including inputs above
the observation scale. We therefore compare whole returns on any
positive-scale compact interval where the return estimates apply.

\begin{lemma}[A whole return and all its prefixes]
\label{path:ordered-return}
Let $K\Subset(0,\delta_{\rm reset})$ lie below the inner return
cutoffs. Consider actual inputs with scales $d\in K$, generated within
the return domains, and with
complete same-base defect
$\kappa=|\xi^r-\Gamma_\lambda(d)|+
|\eta_B-\mathcal B_\lambda(d)|\to0$.  There is an increasing
homeomorphism $h:[0,3]\to[t_0,t_1]$ for their actual return, with
\begin{equation}
 \begin{gathered}
 h(0)=t_0,\quad h(1)=t_{\rm sl},\quad h(2)=t_L,\quad h(3)=t_1,
 \\
 \sup_{\alpha\in[0,3]}
 |\cX_\eps(h(\alpha))-\cC_{\lambda,d}(\alpha)|
                  \le C_K(\eps+r+\kappa).
 \end{gathered}
 \label{path:ordered-estimate}
\end{equation}
Here $t_{\rm sl}$ is the actual first descending $u_r/\eps=M_{\rm sl}$ event,
$t_L$ the actual first release and $t_1$ the next first touchdown.
On compression $h$ is fixed by
\begin{equation}
 J(h(\alpha))=\alpha J_{\rm sl},
 \quad J(t)=\int_{t_0}^t\lambda_N\,\dd t,
 \quad J_{\rm sl}=J(t_{\rm sl}),\qquad 0\le\alpha\le1;
 \label{path:impulse-clock}
\end{equation}
on release and flight it is
\begin{equation}
 h(1+q)=t_{\rm sl}+q(t_L-t_{\rm sl}),\qquad
 h(2+q)=t_L+q(t_1-t_L),\quad 0\le q\le1.
 \label{path:event-fractions}
\end{equation}
The same estimate pairs every stopped prefix in both directions.
Moreover, in the normal coordinates $y,z=\eps d^{3/2}(Y,Z)$,
$v=d^{3/2}V$, at the actual first release,
\begin{equation}
 \left|v_L-d^{3/2}V_d^{\rm n}(R_d)\right|
 +\left|z_L/\eps-d^{3/2}Z_d(R_d)\right|
                         \le C_K(\eps+r+\kappa).
 \label{path:normalized-release}
\end{equation}
The flight durations have a common positive lower bound on $K$.
\end{lemma}

\begin{proof}
Write $\chi=\eps+r+\kappa$.  On the fixed positive compact set,
$e=\eps/d^{3/2}=O_K(\eps)$ and all weighted normal estimates
give unweighted estimates with constants depending on $K$.
By \eqref{path:contact-tube}, the contact lasts $O_K(\eps)$,
and its coefficients have bounded drift.  For $t\le t_{\rm sl}$, put
$f_\theta(w)=1-\Phi_\theta(w)$ and
$\mathcal H_\theta(w)=\int_{M_{\rm sl}}^w f_\theta(q)\,\dd q$.
The profile tail gives $0\le\mathcal H_\theta\le C$.
The fixed slip threshold of Appendix~\ref{sec:contact} gives
$H_\theta(u_r/\eps,X)\ge\kappa_H>0$ on this interval, and
$\dot u_r=a_T-H_\theta(u_r/\eps,X)\lambda_N$.
Since $J(t)\le C_K$,
the identity for $\eps\mathcal H_\theta(u_r/\eps)$ implies on
every prefix
\begin{equation}
 0\le\int_{t_0}^t f_\theta(u_r/\eps)\lambda_N\,\dd t
 \le\kappa_H^{-1}\left\{\int_{t_0}^t f_\theta(u_r/\eps)a_T\,\dd t
             -\eps[\mathcal H_\theta(u_r/\eps)]_{t_0}^t\right\}
 \le C_K\eps.
 \label{path:tail-impulse}
\end{equation}
The two exact velocity rows consequently give, simultaneously on
all compression prefixes,
\begin{equation}
 u_r(t)=u_r(t_0)-h_dJ(t)+O_K(\eps),\qquad
 v(t)=v(t_0)-adJ(t)+O_K(\eps).
 \label{path:impulse-lines}
\end{equation}
Coefficient drifts multiply a bounded impulse, while the free-field
integrals and the unsaturated terms are controlled as above.  At
$t_{\rm sl}$, $u_r=\eps M_{\rm sl}$, so
$J_{\rm sl}=u_r(t_0)/h_d+O_K(\eps)\ge j_K>0$.
The force is positive on $(t_0,t_L)$; hence $J/J_{\rm sl}$ is continuous
and strictly increasing through compression.  This proves that
\eqref{path:impulse-clock} is a homeomorphism and, by
\eqref{path:impulse-lines}, gives precisely the linear impulse
parameter of $\cC_{\lambda,d}|_{[0,1]}$.  Reconstructing
$\omega=(u_r+r-U_*)/\cos\phi$ and using
$|y|+|z|\le C_K\eps$ proves \eqref{path:ordered-estimate} there.

From the actual $M_{\rm sl}$ event use elapsed normal time
$\tau=(t-t_{\rm sl})/\eps$.  The full inherited-state estimate
\eqref{con:M-admission} first compares to the inherited reference
for that actual normalized touchdown input.  Move this reference to the graph
input afterwards.  On $K$, its input sensitivity and
\eqref{con:release-comparison} give
\begin{equation}
 |(Y,V,Z)(\tau)-N_d(\tau)|+
 |L_\eps(\tau)-L_d(\tau)|+
 |L_\eps'(\tau)-L_d'(\tau)|\le C_K\chi
 \label{path:normal-prefix}
\end{equation}
on every stopped release prefix.  The positive-force interval and
strict first-root bracket yield
\[
 R_\eps=(t_L-t_{\rm sl})/\eps,
 \qquad |R_\eps-R_d|\le C_K\chi.
\]
For $1\le\alpha\le2$, compare at actual time
$(\alpha-1)R_\eps$ and then move only the linear reference to
$(\alpha-1)R_d$.  The strip bound gives $u_0=r+O_K(\eps)$,
and the remaining physical coordinates follow from the chart and
\eqref{path:normal-prefix}.  This proves the release portion of
\eqref{path:ordered-estimate}, including its moving endpoint,
and also \eqref{path:normalized-release}.  In particular
$|z_L|\le C_K\eps$.  The opening margins give
$\min(v_L,v_L+z_L/\eps)\ge v_K>0$.

On the analytic free extensions the mechanical field is smooth and
independent of belt translation.  Gronwall, on their common compact
time interval, compares the actual-release and graph-release flows
with error $C_K\chi$.  Their exact effective gap is
\begin{equation}
 g(\ell)=z_L(1-e^{-\ell/\eps})+
                  \int_0^\ell v(q)\,\dd q.
 \label{path:exact-free-gap}
\end{equation}
To exclude an additional root next to release, use $|b|\le B_K$:
\[
 g'(\ell)\ge\min(v_L,v_L+z_L/\eps)-B_K\ell>0
\]
on a fixed short opening interval.  Thereafter the reference gap
has a positive minimum away from its terminal root.  In a common
terminal bracket its velocity has a negative margin, and the memory
derivative is uniformly negligible because that bracket is bounded
away from release.  The actual gap has opposite endpoint signs
and negative derivative there.  Its first subsequent zero therefore
is the actual touchdown, with
$|T_\eps-T_d|\le C_K\chi$.  This also gives a positive common
lower bound for $T_\eps$.  Moving the reference time from
$qT_\eps$ to $qT_d$ proves the flight portion of
\eqref{path:ordered-estimate}.  The three clocks agree at their
joins and are strictly increasing.  Restricting the same pairing
to any prefix proves both image inclusions as well as its ordered
comparison.
\end{proof}

Apply this lemma to the last input in \eqref{path:last-input} and
write $\alpha_\eps=h_\eps^{-1}(t_s)$.  Compactness gives
subsequences $(d_\eps,\alpha_\eps)\to(d,\alpha)$, and the actual
section equation gives
\begin{equation}
 p(\cC_{\lambda,d}(\alpha))=-as.
 \label{path:actual-stop-equation}
\end{equation}
The contact has $p=-ad$, and the free part decreases strictly from
$-ad$ to $-aF_\lambda(d)$.  Its complete solution set is therefore
\begin{equation}
 \mathcal A_s=
 \{(d,j_s(d)):F_\lambda^{-1}(s)\le d\le s\}
                   \cup(\{s\}\times[0,2]).
 \label{path:terminal-domain}
\end{equation}
This includes a crossing at the next touchdown even if the varying
reference curve has not itself crossed before its own endpoint.
Every terminal prefix converges in order, since
\begin{equation}
 \sup_{q\in[0,1]}
 |\cX_\eps(h_\eps(q\alpha_\eps))
                       -\cC_{\lambda,d}(q\alpha)|\longrightarrow0.
 \label{path:terminal-prefix}
\end{equation}
For $\alpha=0$ the limiting piece is constant, as allowed by the
completed ordered topology.

\subsection{Actual realization of the terminal domain}

Reduce $s_0$ once more so that $2s_0$ lies below all inner input
cutoffs and, for $0<s\le s_0$,
\begin{equation}
 2s<\Delta,\qquad (2s/c_D)^{4/3}<s/3,\qquad
 c_D(s/2)^{3/4}>s.
 \label{path:scan-margins}
\end{equation}

\begin{lemma}[A continuous sweep of one actual contact]
\label{path:continuous-sweep}
For each fixed complete label and every sufficiently large $j$
there is an interval $J_j=[a_j,c_j]$, $c_j\to0$, on which the same
indexed actual touchdown, contact and next return exist continuously,
and
\begin{equation}
 \delta_j(a_j)=s/2,\qquad\delta_j(c_j)=2s,
 \qquad s/2\le\delta_j(\eps)\le2s\quad(\eps\in J_j).
 \label{path:scan-band}
\end{equation}
Every earlier controlled touchdown is below $s/3$, and the next
touchdown is above $s$.  The ordered comparison
\eqref{path:ordered-estimate} holds uniformly on $J_j$ with error
tending to zero as $j\to\infty$.
\end{lemma}

\begin{proof}
At positive compliance the positive-part vector field is locally
Lipschitz and depends continuously on compliance.  The prescribed
initial data are continuous.  The strict opening, closing, first
$M_{\rm sl}$ and first-release brackets consequently give continuous
dependence of any fixed finite sequence of the actual events.
The needed local constants are only on sets bounded away from zero
in compliance and input scale.  Positive gap between successive
first events prevents a change of the indexed passage.

Fix a common $\eps_*>0$ so that the original entrance and its
admission hold for $0<\eps<\eps_*$. For fixed $j$, the entrance
and growth upper bound imply $\delta_j(\eps)\to0$; all preceding
inputs eventually satisfy the return hypotheses. Let $O_j\subset(0,\eps_*)$ be the
open set on which the events through $\TD_j$ are generated with
$\delta_n<\Delta$ for $3\le n<j$. Its component adjacent to zero
is $(0,b_j)$. These components are nested. For every fixed
$\eps_0\in(0,\eps_*)$, finite exit at that compliance excludes
$\eps_0$ from $O_j$ for sufficiently large $j$; hence $b_j\to0$.
In particular $b_j<\eps_*$ for all sufficiently large $j$, so its
right endpoint is not the artificial end of the total compliance
interval. At such a positive endpoint $b_j$, the finitely many scales
have a positive lower bound and their full states and strict margins
stay in compact sets.  Extend the finite events to that endpoint
successively.  Inputs equal to $\Delta$ are allowed because
$\Delta<\delta_{\rm reset}$.  No earlier input
$\delta_n(b_j)$, $n\le j-2$, can equal $\Delta$: its next output
would be larger than $\Delta$, contradicting the left limiting
inequality for that next input.  If also
$\delta_{j-1}(b_j)<\Delta$, all finite events and strict input
inequalities would continue past $b_j$.  Hence
\[
 \delta_{j-1}(b_j)=\Delta,
 \qquad \delta_j(b_j)>\Delta>2s.
\]
Only this endpoint call uses $\Delta$; its outer output is not
used as a new input.

Take the first compliance $c_j<b_j$ at which $\delta_j=2s$ and
the last preceding compliance $a_j$ at which $\delta_j=s/2$.
These first and last choices give \eqref{path:scan-band} on the
whole interval by continuity. Its scanned input is
itself below $\Delta$, so the return estimates apply to the next return throughout.
The growth lower bound and \eqref{path:scan-margins} give
\[
 \delta_{j-1}\le(2s/c_D)^{4/3}<s/3,
 \qquad\delta_{j+1}\ge c_D(s/2)^{3/4}>s.
\]
Earlier scales are smaller by strict growth.  Finally the budget
gives $B_j\le q_{\rm g}^{j-3}B_3+C\Omega_\eps\to0$ uniformly on
$J_j$; \eqref{path:dependent-row} recovers the full defect on
$[s/2,2s]$.  Lemma~\ref{path:ordered-return} applies on this entire
positive compact set, proving the last assertion.
\end{proof}

\begin{proposition}[Every terminal pair is realized]
\label{path:realization}
For every fixed complete label,
\[
 \Clust_{\eps\downarrow0}(d_\eps,\alpha_\eps)=\mathcal A_s.
\]
For fixed $\lambda,s$ and $d\in(d_*,s)$, $d_*=F_\lambda^{-1}(s)$,
the explicit sequence
\begin{equation}
 \eps_m(d)=\exp\{-\tfrac{32}{9}(3/4)^{-m}
                                  \mathcal S_\lambda(d)\}
 \label{path:common-free-sequence}
\end{equation}
realizes $(d,j_s(d))$ uniformly on the fixed compact
$(k,\theta,\nu,\mathbf n)$ class. Contact and endpoint
realization is asserted for each fixed complete label and target;
Proposition~\ref{sel:threshold-selection} supplies explicit common
sequences for the two free endpoints.
\end{proposition}

\begin{proof}
Fix first $\alpha\in(0,2]$.  On $J_j$ define $t_{j,\alpha}$
by \eqref{path:impulse-clock} if $0<\alpha\le1$, and by the
first formula of \eqref{path:event-fractions} if $1\le\alpha\le2$.
These actual event times are continuous in compliance.  In the
first case this follows from continuity of the flow and integral,
and the strictly increasing inverse of $J/J_{\rm sl}$ with $J_{\rm sl}\ge j_s>0$.
The two definitions agree at $\alpha=1$.  The continuous function
\[
 G_j(\eps)=p(\cX_\eps(t_{j,\alpha}(\eps)))+as
\]
has signs
\begin{equation}
 G_j(a_j)=as/2+O_s(a_j)>0,\qquad
 G_j(c_j)=-as+O_s(c_j)<0
 \label{path:phase-scan-signs}
\end{equation}
for large $j$.  Indeed contact takes $O_s(\eps)$ and $\dot p$ is
bounded throughout its actual tube.  Take a root $\eps_j\in J_j$.
Since $\alpha>0$, $t_{j,\alpha}>t_j$; strict decrease of $p$ gives
$\delta_j<s$.  The predecessor and successor bounds in the sweep
exclude both neighboring last-input indices: $N_s=j+1$, and the
root is the first physical section time.  Moreover
$\delta_j=s+O_s(\eps_j)\to s$ and the stopping phase equals
$\alpha$ exactly.  The comparison lemma proves its whole terminal
prefix limit.  By passing to later indices the compliances can be
made strictly decreasing.

For $\alpha_i\downarrow0$, choose sufficiently late terms from
these positive-target sequences with decreasing compliance and
all base and comparison errors below $1/i$.  This realizes $(s,0)$.
Solving the equation directly at $\alpha=0$ would instead give
$\delta_j=s$ and $N_s=j$, with the predecessor as last input;
the positive-target diagonal avoids that change of index.

For the free targets set $j=m+3$ at
\eqref{path:common-free-sequence}. The exact graph chain started at
the actual $\delta_3$ has phase at $j$ tending uniformly to
$\mathcal S_\lambda(d)$, by \eqref{sel:uniform-entry-phase}.
Its bases through $j$ are eventually below $3s/2$.
Lemma~\ref{sel:sharp-phase} therefore gives the actual events through
this index and
\begin{equation}
 \begin{gathered}
 \mathcal S_\lambda(\delta_j)
 =\mathcal S_\lambda(d)+\rho_c^m
       \{\chi_\lambda(\mathfrak a)+r_{\mathfrak a}(\eps_m)\}
             +O(\eps_m^{3/32})\longrightarrow\mathcal S_\lambda(d),\\
 E_j=O(\eps_m^{3/32}).
 \end{gathered}
 \label{path:common-free-budget}
\end{equation}
Thus $\delta_j\to d$ and its complete independent state tends to
the graph, uniformly. Equation~\eqref{path:dependent-row} recovers
the dependent acceleration row. The preceding graph base tends to
$F_\lambda^{-1}(d)<d_*$, so the preceding actual inputs stay a
fixed distance below $s$.

This output is separately checked as a new input: it is close to
$d<s$, below the reset cutoff, retains the full cone and closing
margins, and satisfies the propagated compliance and detuning
inequalities.  Its actual next return has output tending to
$F(d)>s$.  The preceding scales are below $d_*<s$, so
$N_s=m+4$ and $d_\eps\to d$.  Every limit of its stopping phase
satisfies \eqref{path:actual-stop-equation}.  Since $d<s$, the
unique solution is the free parameter $j_s(d)$.  This proves
uniform convergence: a failure of uniformity would give a varying
compact-label subsequence and a different phase limit, contrary to
the same equation and unique free root.  The complete tuple and
ordered prefix comparisons are uniform for the same reason.
For either free endpoint, a diagonal of interior targets and
successively late terms realizes that endpoint with the complete
label fixed.

We have realized the entire vertical contact segment and
free branch in \eqref{path:terminal-domain}.  The reverse inclusion
was already proved by \eqref{path:actual-stop-equation}; this proves
the stated equality of cluster sets.
\end{proof}

\subsection{Backward matching and the whole stopped path}

\begin{lemma}[Every fixed number of last returns]
\label{path:backward-match}
If $(d_\eps,\alpha_\eps)\to(d,\alpha)$, then for every fixed
$M$ and $0\le m\le M$,
\begin{equation}
 \delta_{N_s-1-m}\longrightarrow b_m=F_\lambda^{-m}(d).
 \label{path:backward-bases}
\end{equation}
The complete normalized touchdown state tends to the graph at each $b_m$.
The $M$ corresponding full actual returns converge in order to
$\cC_{\lambda,b_m}$, and the last prefix converges to
$\cC_{\lambda,d}|_{[0,\alpha]}$.
\end{lemma}

\begin{proof}
The last scale is bounded below by $c_s$.  Repeated use of the
growth upper bound gives a positive lower bound for each fixed
finite number of predecessors, depending on $M,s$.  Their indices
tend to infinity.  The graph budget and the dependent-row recovery
therefore give complete graph attraction at each such input.
Given the limit $b_{m-1}$, every limit $b$ of its predecessor
satisfies $F_\lambda(b)=b_{m-1}$ by
\eqref{path:same-base-output}.  Strict increase of $F_\lambda$
makes this limit unique, equal to $F_\lambda^{-1}(b_{m-1})$.
Compactness removes further extraction and proves
\eqref{path:backward-bases}.  Lemma~\ref{path:ordered-return}
applies to these finitely many positive-scale compact intervals;
continuity in their bases gives the fixed reference returns.
Equation~\eqref{path:terminal-prefix} treats the last prefix.
\end{proof}

\begin{proof}[Proof of Theorems~\ref{thm:whole} and~\ref{thm:common-free}]
Lemma~\ref{path:first-section} gives the first transverse section,
and Lemma~\ref{path:stopped-cascade} gives $N_s\to\infty$.
Take a sequence with terminal pair tending to $(d,\alpha)$.
The pair lies in $\mathcal A_s$ by
\eqref{path:actual-stop-equation}.  Fix a depth $M$, and put
$b_M=F^{-M}(d)$.  For sufficiently small compliance all actual
inputs earlier than $\TD_{N_s-1-M}$ are at most $2b_M$.
Their whole return segments are therefore in the same ball
$G+B_{C(2b_M)^{3/4}}$, by
\eqref{path:contact-tube}--\eqref{path:free-normal-tube}.
This is a supremum over all early returns.  The analogous infinite
reference tail in $L_d$ is in $G+B_{Cb_M^{3/4}}$.

For a fixed incoming cutoff $\eta>0$,
Lemma~\ref{path:first-section} aligns the incoming prefixes and puts
the remaining original-to-$\TD_3$ curve in
$G+B_{C\eta+o(1)}$.  Match in their order the entire intervening
actual curve and the remaining incoming reference tail followed
by the small infinite reference cascade.  Their common ball gives
an error at most $C\eta+Cb_M^{3/4}+o(1)$, regardless of how many
actual returns it contains.  Match the remaining $M$ full returns
and final prefix by Lemma~\ref{path:backward-match}, and join
their increasing pairings at the event endpoints.  Their finite
maximum error tends to zero.  Constant limit pieces are handled by
nondecreasing continuous surjections, approximated by increasing
homeomorphisms as in the definition of the completed distance.
Consequently
\begin{equation}
 \limsup_{\eps\to0}\dF(f_\eps,
                  \cW_{\lambda,\nu,s}(d,\alpha))
                         \le C\eta+Cb_M^{3/4}.
 \label{path:whole-error}
\end{equation}
First let $\eps\to0$ at fixed $M,\eta$, and then let
$\eta\downarrow0$, $M\to\infty$.  Inverse contraction gives
$b_M\to0$, uniformly for $d\in[d_*,s]$ and the mechanical compact
class.  This proves the whole ordered convergence.  Since
$d_H\le\dF$, it also proves the two-sided Hausdorff convergence.

The map $(d,\alpha)\mapsto\cW_{\lambda,\nu,s}(d,\alpha)$ is
continuous on $\mathcal A_s$: for fixed $M$ use continuity of the
finitely many last reference pieces and their terminal prefix,
then bound both omitted tails uniformly by
$C\sup_{d\in[d_*,s]}(F^{-M}(d))^{3/4}$.  Its image is compact.
Every sequence tending to zero has a convergent terminal-pair
subsequence, giving the upper cluster inclusion.  At a positive
compliance limit, continuous dependence and the transverse section
give continuity of the stopped curve.  These two cases prove
relative compactness of the sufficiently small compliance family.
Proposition~\ref{path:realization} gives the reverse inclusion from
actual compliances with the same complete label.  This proves
the exact cluster classification for paths and for their images.

For the common free sequence, the entrance estimate, the budgets,
the incoming comparison, the fixed-depth comparison and both small
tail bounds are uniform on $(k,\theta,\nu,\mathbf n)$ in the fixed
compact class.  The incoming prefix is always compared with its
own $I_{\lambda,\nu}$.  Thus \eqref{path:whole-error} gives the
asserted uniform whole-path convergence. The same finite-depth
comparison and small-tail bounds apply whenever the terminal pairs
converge uniformly on the compact auxiliary class. We use this
observation for the explicit endpoint sequences in
Proposition~\ref{sel:threshold-selection}. The free sequence fixes
$\lambda,s,d$; the contact roots above are taken for each fixed
complete label and target.
\end{proof}



\begin{corollary}[Different limits from the same full initial state]
\label{cor:fixed-initial-paths}
Fix $\lambda,k,\theta,\nu$ and take $\mathbf n=0$. The common full
initial state is $(\phi_\nu,0,\omega_\nu,0,0)$. There are two
compliance sequences along which the stopped paths tend respectively
to $\cW_{\lambda,\nu,s}(s,0)$ and
$\cW_{\lambda,\nu,s}(s,1)$. Their terminal velocities and compact
path images differ.
\end{corollary}

Here the belt still follows $U_b=U_*-k\sqrt\eps$.

\begin{proof}[Proof of Corollary~\ref{cor:fixed-initial-paths}]
With $\mathbf n=0$, \eqref{eq:bd-physical-initial-data} is the stated
fixed full initial point. Theorem~\ref{thm:whole} realizes the two
targets with the remaining label fixed. Put
$U_s=d_\lambda s(1+s^{1/3}\Gamma_U(s))>0$ and
$h_s=\mu A(\phi_s)-B(\phi_s)>0$, where $p(\phi_s)=-as$.
The compression curve in \eqref{gr:compression} gives
\[
 \omega_{(s,0)}-\omega_{(s,1)}=\frac{U_s}{\cos\phi_s}>0,
 \qquad
 v_{(s,0)}-v_{(s,1)}=\frac{as}{h_s}U_s>0.
\]
The first image meets $p=-as$ only at its admission point; the second
contains the different arrest point on that section. Hence their
compact path images differ.
\end{proof}
